% Reconstructed text from V5 PDF page 19; non-native source
\par\smallskip\noindent\rule{\textwidth}{0.4pt}\par{\sffamily\footnotesize Original V5 page 19 -- recovered text}\par\begin{flushleft}\footnotesize\ttfamily
\par\smallskip
\par\smallskip
\noindent Diagnostics and Serving Infrastructure for MAS 19\par
\par\smallskip
\par\smallskip
\par\smallskip
\noindent Tooling. AgentTrace [2] implements three logging surfaces — operational (method-level ex-\par
\noindent ecution tracing with paired start/complete events), cognitive (LLM reasoning chain extrac-\par
\noindent tion including CoT parsing, confidence estimates, and plan/reflection fields), and contextual\par
\noindent (outbound API interactions auto-instrumented via OTel monkey-patching) — unified under a\par
\noindent schema-validated append-only envelope with UUID, surface type, trace/span IDs, and UTC\par
\noindent timestamps (§3.1–3.6). The system produces dual output: JSONL files for local analysis and\par
\noindent OTel spans for distributed backends. While the architecture is well-specified with pseudocode\par
\noindent (Algorithm 1, §4.2), no empirical evaluation of overhead, latency, or storage costs is reported\par
\noindent (§4.5 states only “low overhead” qualitatively). No repository link or package name is provided.\par
\noindent  Insights Generator [23] provides analytics over collected execution traces, enabling aggre-\par
\noindent gate analysis across multiple agent runs: failure frequency by step type, time-to-failure dis-\par
\noindent tributions, and cross-agent correlation patterns. The view-oriented debugging compiler [46]\par
\noindent transforms raw execution logs into interactive visual representations: the compiler applies view\par
\noindent transformations (filtering, aggregation, temporal alignment) to produce domain-specific debug-\par
\noindent ging views from a common trace format, enabling human inspectors to navigate complex multi-\par
\noindent agent executions without manually parsing log files.\par
\noindent  LangSmith (LangChain, commercial tooling) provides a production-grade trace collection\par
\noindent and analysis platform with built-in LLM evaluation, but its proprietary nature limits independent\par
\noindent verification of overhead and accuracy claims.\par
\noindent  The observability layer’s principal limitation for cross-layer analysis is coverage: no ex-\par
\noindent isting standard captures the serving-side state (KV cache residency, eviction events, memory\par
\noindent pressure) that Section 6 identifies as relevant for diagnostic-infrastructure interaction. Extend-\par
\noindent ing OTel GenAI conventions to include serving-side attributes is one of the open directions\par
\noindent proposed in Section 7.\par
\par\smallskip
\noindent 4.3 Safety Intervention and Guardrails\par
\noindent Safety systems actively monitor and intervene in agent execution to prevent harmful actions.\par
\noindent We review eight Tier-1 systems spanning three architectural patterns: graph-based guardrails,\par
\noindent modular safety agents, and control-theoretic recovery.\par
\par\smallskip
\noindent Graph-based guardrails. G-Safeguard [37] addresses a specific threat model: adversarial\par
\noindent agents embedded within a multi-agent system that manipulate other agents through malicious\par
\noindent utterances. The system constructs a multi-agent utterance graph from inter-agent communica-\par
\noindent tion, using SentenceBERT/MiniLM embeddings as node features and a learnable permutation-\par
\noindent invariant fusion of interaction histories as edge features. An edge-featured GNN with a soft\par
\noindent probabilistic classifier detects attacker nodes after each dialogue round. Trained on 800 samples\par
\noindent from a GPT-4o-mini-based 8-agent MAS, the system scales to 20–80 agents without retraining.\par
\noindent A key limitation is that detection is reactive: “by the time G-Safeguard identifies an attacker,\par
\noindent certain nodes within the MAS have already been successfully compromised.” The system oper-\par
\noindent ates purely on the utterance graph and does not inspect internal model serving state (ACL 2025\par
\noindent Long Papers).\par
\par\smallskip
\noindent Modular safety agents. ShieldAgent [8] represents safety policies as Linear Temporal Logic\par
\noindent (LTL) formulas organized into action-based probabilistic rule circuits via Markov Logic Net-\par
\noindent works. For each agent action, the system extracts action/state predicates from policy documents,\par
\noindent applies spectral clustering on state predicates for rule grouping, and evaluates via two fine-\par
\noindent tuned guardrail models on InternVL2-2B. The Autonomous Safety Policy Manager (ASPM)\par
\end{flushleft}
